\chapter{Dopant Transferability and Validation Program}
\index{transferability!dopant|textbf}

\label{ch:dopant-transferability}

\chapterstatus{proposed phosphorus demonstration and boron transferability
program; no production impurity reduction or validation result is claimed.}

\section{Phosphorus as the first impurity demonstration}
\index{phosphorus!demonstration}


After the bulk representation and alignment gates pass, phosphorus is the first
complete impurity path. The branch requires a compatible P pseudopotential,
neutral spin-polarized P$_\mathrm{Si}^{0}$ supercells, controlled internal
relaxation, matched pristine comparison cells, and supercell-size evidence. The
aligned difference is then decomposed into nested scalar, orbital, onsite,
hopping, and range-restricted model classes.

The objective is not to assume that a screened Coulomb potential plus one
central-cell parameter is sufficient. It is to determine which operator and
bound-state information each tested class preserves. A continuum donor model
becomes a first-principles reduction only after its host parameters, embedding,
exterior error, cross-coupling error, and bound-state comparison are tied to the
accepted atomistic parent.

\section{Boron as a transferability branch}
\index{boron!transferability}


Boron tests whether the methods transfer to a physically different band-edge
problem. The final acceptor branch requires a fully relativistic noncollinear
spinor treatment with SOC and a compatible valence subspace. The phosphorus
basis or scalar model class cannot be transferred unchanged merely for
convenience.

Comparing the P and B hierarchies may reveal which reduction features are common
to substitutional impurities and which are specific to conduction- or
valence-manifold structure. Such a comparison becomes eligible only after each
branch has its own accepted atomistic operator and minimal-model evidence.

\section{Methodological implications}

\subsection{Structured learning as a later diagnostic}

After an aligned impurity operator exists, structured learning may provide an
expressive comparison class for the reduction hierarchy. The proposed network
does not solve the Kohn--Sham parent problem and does not replace the validated
identification map. It predicts parameters of a deterministic Hamiltonian
constructor or a constrained residual beyond a simpler physical model.

The relevant result would be the sequence of attainable errors across frozen
scalar, orbital, finite-range, continuum, and learned classes, together with
separate optimization and withheld-evaluation evidence. A learned model that
improves on effective-mass theory near the impurity but offers no material
advantage in an accepted exterior domain could contribute to the crossover
analysis. Such a pattern is proposed work, not a result of the present draft.
Joint learning of the gauge and reduced operator is deferred until the
fixed-identification problem is understood, because low joint training loss
would not by itself establish a physically meaningful correspondence.

\subsection{Composition and path dependence}
\index{composition!reduction paths}


A completed hierarchy still requires cross-path assessment. Direct and
Wannier-mediated bulk fits may disagree. Extracting an impurity before reduction
may differ from reducing pristine and doped parents before subtraction.
Successive lattice and continuum reductions may amplify earlier errors.

These noncommuting diagrams are not expected to vanish by definition. Their
paths, aligned outputs, norms, and validation domains must be measured. The
final compositional claim may be that a diagram commutes within tolerance over a
declared domain, that it fails by a quantified defect, or that available
metadata are insufficient to compare the paths.

\subsection{Negative results as supported outcomes}
\index{negative result}


The research program is designed to retain negative findings. An unacceptable
nearest-neighbor model, an unstable alignment, a nonconverged impurity sequence,
an empty admissible-set intersection, or an infinite crossover radius over the
tested domain can all be valid outcomes. Their value depends on preserving the
frozen model hierarchy, rejected candidates, tested range, metrics, and
limitations.

A negative finding is not generalized beyond its tested classes and domain. A
failed software contract is corrected; a failed numerical tolerance is refined
or reported; a physically inadequate model class is not repaired by weakening
its acceptance rule after the result is known.
